% Inputtable proof fragment; standalone wrapper: excess-main.tex.
\section*{Scope of this note}
We independently derive the proposed replacement of the extremal ratio by
an additive excess over a fixed target. Attainment and closed-support
contact hold for every positive target at each fixed index. The conserved
identity has the coefficient $(n+H^2)a$, with $H$ the raw sum threshold.
This removes the actual score from the contact and remainder coefficients.
It does not establish a new uniform local exclusion or transfer the old
numerical contracts to a different maximizing law. All original source
files and certificates are preserved.

\section{Definitions and exact scope}
Let $\mathcal P$ be the laws on $\R$ with $\E X=0$, $\E X^2=1$ and
$\rho=\E|X|^3<\infty$. All sums below use independent copies of one
marginal law. For $S_n=X_1+\cdots+X_n$, set
\[
 F_{n,P}(H)=\Pp(S_n\le H),\qquad z=H/\sqrt n,
 \quad R_n(P,H)=\frac{\sqrt n}{\rho}\{F_{n,P}(H)-\Phi(z)\}.
\]
The global Berry--Esseen constant at this index is the supremum of the
absolute standardized discrepancies, divided by $\rho/\sqrt n$; call it
$C_n$. Reflection and approximation by closed thresholds show that it is
also $\sup_{P,H}R_n(P,H)$. For a fixed number $c>0$ define
\[
 a(P,H;c)=\sqrt n\{F_{n,P}(H)-\Phi(z)\}-c\rho,
 \qquad A_n(c)=\sup_{P,H}a(P,H;c).
\]
The excess $a$ has the units of $\rho$ times a score difference. It is
\emph{not} the existing manuscript's $a_n=C_n-\CE$. At a maximizing pair
of this new problem its actual positive score is
\begin{equation}\label{eq:actual}
 C_{\rm act}=c+a/\rho.
\end{equation}
In general $C_{\rm act}$ need not equal $C_n$.

For conservation and defect positivity we later specialize to
$c=\CE=(3+\sqrt{10})\phi(0)/6$. Attainment and contact themselves require
only $c>0$. No Fourier bound, numerical acceptance or uniform index limit
is part of their proof.

\section{Attainment and the equality boundary}
\begin{theorem}\label{thm:attain}
For each fixed integer $n\ge1$ and $c>0$, the maximum $A_n(c)$ is
attained by a pair $(P,H)$ with a positive closed-CDF discrepancy.
Moreover,
\[
 A_n(c)>0\iff C_n>c,\qquad
 A_n(c)=0\iff C_n=c,\qquad A_n(c)<0\iff C_n<c.
\]
\end{theorem}
\begin{proof}
Lyapunov gives $\rho\ge1$, hence $A_n(c)\le\sqrt n-c$. Symmetric
Bernoulli, at $H=n$, gives the stronger lower bound
\[
 A_n(c)\ge\sqrt n\{1-\Phi(\sqrt n)\}-c>-c.
\]
Write $A=A_n(c)$, and choose $a_j=a(P_j,H_j;c)\to A$. For all sufficiently
large $j$, $a_j\ge A-1$ and
\[
 c\rho_j\le\sqrt n-a_j\le\sqrt n-A+1,
 \qquad \sqrt n\{F_{n,P_j}(H_j)-\Phi(H_j/\sqrt n)\}
       =a_j+c\rho_j\ge a_j+c\ge\eta>0,
\]
where $\eta$ can be any fixed number smaller than $A+c>0$ once $j$ is
large enough. Unit variance implies tightness. A bound $\rho_j\le M$
gives
\[
 \int_{|x|>K}x^2\,dP_j\le M/K,
 \qquad\int_{|x|>K}|x|\,dP_j\le M/K^2.
\]
Take subsequences $P_j\Rightarrow P$ and $\rho_j\to b$. Bounded continuous
truncation and these tails preserve mean zero and variance one;
lower semicontinuity gives $1\le\rho(P)\le b<\infty$.

For $H_j<0$, Chebyshev gives $F_{n,P_j}(H_j)\le n/H_j^2$.
For $H_j>0$, the positive discrepancy is at most
$1-\Phi(H_j/\sqrt n)\le n/H_j^2$. Its lower bound $\eta/\sqrt n$ thus
bounds $|H_j|$. Take $H_j\to H$. Fixed $n$ permits passage from marginal
weak convergence to the sum law. For every $t>0$, eventually
$H_j\le H+t$, so the closed-set inequality and then right continuity give
\[
 \limsup_j F_{n,P_j}(H_j)\le F_{n,P}(H+t)
       \longrightarrow F_{n,P}(H).
\]
Consequently
\[
 a(P,H;c)\ge A+c\{b-\rho(P)\}\ge A.
\]
The definition of $A$ forces equality, and $c>0$ forces $\rho(P)=b$.
Third-moment convergence is a \emph{conclusion of maximality}; bounded
third moments were not used to assert uniform integrability of third
powers. Also $A+c\rho\ge A+c>0$, proving positive discrepancy.

Reflection changes a negative closed discrepancy into a positive open
discrepancy of the reflected law. Closed thresholds increasing to that
open threshold approximate it; the penalty is unchanged. This proves
the stated supremum identity for $C_n$ without presuming continuity at an
atom. If $C_n<c$, every excess is at most
$\rho(C_n-c)\le C_n-c<0$. If $C_n>c$, a positive closed score above $c$
gives positive excess. If $C_n=c$, choose scores $r_j\uparrow c$ with
$r_j\ge c/2$. Their moments satisfy $\rho_j\le2\sqrt n/c$, and
$\rho_j(r_j-c)\to0$. Thus $A_n(c)\ge0$, while all excesses are at most
zero. These observations prove all three equivalences; no ratio-attainment
lemma is needed for this last step.
\end{proof}

\section{Atom-safe fixed-target contact}
Fix a maximizing pair of Theorem~\ref{thm:attain}, with $n\ge2$. Set
$F=F_{n,P}(H)$, $\varphi=\phi(z)$,
$\beta=\E[X|X|]$ and $\gamma=\E X^3$.
\begin{theorem}\label{thm:contact}
The following function is nonpositive for every $y\in\R$ and vanishes
on the entire closed support of $P$:
\begin{align}
 I_c(y)={}&n^{3/2}\{F_{n-1,P}(H-y)-F\}+ny\varphi
       +\tfrac12 H\varphi(y^2-1)\notag\\
 &+\tfrac32c\rho(y^2-1)-c\{|y|^3-\rho-3y\beta\}.
 \label{eq:contact}
\end{align}
That support is compact. Equivalently, for $h(y)=F_{n-1,P}(H-y)$,
$h\le Q_c$ everywhere and $h=Q_c$ on the closed support, where
\begin{align}
 nQ_c(y)={}&nF+\frac{H\varphi+c\rho}{2\sqrt n}
  -\left(\sqrt n\varphi+\frac{3c\beta}{\sqrt n}\right)y\notag\\
 &-\frac{H\varphi+3c\rho}{2\sqrt n}y^2
  +\frac{c}{\sqrt n}|y|^3.\label{eq:Q}
\end{align}
\end{theorem}
\begin{proof}
Contaminate the raw law by
$P_\varepsilon=(1-\varepsilon)P+\varepsilon\delta_y$. Its mean and variance
are $m_\varepsilon=\varepsilon y$ and
$v_\varepsilon=1+\varepsilon(y^2-1)-\varepsilon^2y^2$.
Let $r_\varepsilon=\E_{P_\varepsilon}|X-m_\varepsilon|^3$. The standardized
law, evaluated at the standardized image of the \emph{fixed raw event},
has excess
\[
 G_c(\varepsilon)=\sqrt n\left\{F_{n,P_\varepsilon}(H)
 -\Phi\!\left(\frac{H-nm_\varepsilon}{\sqrt{nv_\varepsilon}}\right)\right\}
 -c\frac{r_\varepsilon}{v_\varepsilon^{3/2}}.
\]
For small $\varepsilon\ge0$ this is admissible; maximality gives
$G_c(\varepsilon)\le G_c(0)=a$. The raw CDF is a polynomial in
$\varepsilon$ and has derivative $n(h(y)-F)$. Dominated differentiation
of $|x-\varepsilon y|^3$, bounded by
$3|y|(|x|+|y|)^2$, gives
\[
 r'_0=|y|^3-\rho-3y\beta,
 \quad \left(\frac{H-nm_\varepsilon}{\sqrt{nv_\varepsilon}}\right)'_0
 =-\sqrt n\,y-\frac{H}{2\sqrt n}(y^2-1).
\]
Differentiation now gives $G'_c(0+)=I_c(y)\le0$. There is no division by
$\rho$ in this fixed-target derivative.

Integrating \eqref{eq:contact} against $P$ gives zero. All its terms are
integrable under the third-moment assumption. The map $h$ is upper
semicontinuous: at a jump the closed event chooses its larger value.
Thus $I_c$ is upper semicontinuous. If it were strictly negative at a
closed-support point, a neighborhood of that point would have positive
$P$-mass and a strictly negative upper bound for $I_c$, contradicting its
zero integral. Hence it vanishes at every such point, including atoms and
accumulation points. Because $c>0$, a bound of the form
$I_c(y)\le-c|y|^3+M_2y^2+M_1|y|+M_0$ excludes both sufficiently distant
rays from the support. Rearrangement gives \eqref{eq:Q}.
\end{proof}

Neither proof assumes $H\ge0$. Reflecting a positive closed discrepancy
again does not in general preserve that event. The actual orientation
must be retained after choosing the maximizing pair.

\section{Predecessor support: the extra term has a favorable sign}
\begin{theorem}\label{thm:support}
Suppose additionally $C_{n-1}\le c$. Put $v=n-1$,
$s=\sqrt{n/v}$ and $e_n=n(s-1)-1/2$. For every support point $y$,
\begin{equation}\label{eq:support}
 0\le\mathcal H_n(H,y)
       +c\{-|y|^3+\tfrac32\rho y^2+3\beta y+\rho e_n\}-na,
\end{equation}
where, with $F_v(x)=\Phi(x/\sqrt v)$,
\[
 \mathcal H_n(H,y)/n^{3/2}
 =F_v(H-y)-F_{v+1}(H)+yF'_{v+1}(H)
                    -\tfrac12(y^2-1)F''_{v+1}(H).
\]
\end{theorem}
\begin{proof}
For the same marginal,
$F_{n-1,P}(H-y)\le F_v(H-y)+c\rho/\sqrt v$.
Substitute this upper bound and
$F=\Phi(H/\sqrt n)+(c\rho+a)/\sqrt n$ in the zero contact equality.
The constant moment contribution is
$c\rho\{n(s-1)-1/2\}$, and the new excess contributes exactly $-na$.
The Gaussian derivatives give the displayed expression. This proves the
claim for each sign of $y$ and every real $H$.
\end{proof}

If $a\ge0$, dropping $-na$ gives the old zero-loss support inequality
with coefficient $c$. Therefore the elementary support bounds in current
Section~4 apply after their hypotheses are restated as
Theorem~\ref{thm:contact}, $C_{n-1}\le c$, $a\ge0$, and the same moment
caps. At $c=\CE$, the lower numerical multiplier $0.4095<c$ remains
valid only after its cubic coefficient is negative, as in the existing
proof. The all-$H$ fourth-order expansion, both ray signs, and frozen
moment caps remain necessary. The structural moment bounds are universal
bounds for each law's actual score, so \eqref{eq:actual} with $a\ge0$
supplies their lower-score premise. This transfer does not re-prove those
external structural inequalities or acquire the later conditional moment
caps automatically.

\section{Independent conservation calculation}
This section needs only the explicit contact relations in
Theorem~\ref{thm:contact}, not the identity of the actual score with the
contact coefficient. Set $\nu=\mathcal L(H-S_{n-1})$ and
$\tau(t)=\E[X\ind_{\{X\ge t\}}]$. Centering makes $\tau\ge0$, and
signwise Fubini gives $\int\tau=1$ and
\[
 \int t\tau(t)\,dt=\gamma/2,\quad
 \int t^2\tau(t)\,dt=M_4/3,\quad
 \int |t|^3\tau(t)\,dt=M_5/4.

\]
Here $M_4=\E X^4$ and $M_5=\E|X|^5$.
These moments are available because compact support has already been
proved. For each bounded complementary gap $(u,v)$ of the closed support,
$\tau$ is constant, say $\tau_{u,v}$. Define
\begin{align*}
 J&=\sum_{(u,v)}\tau_{u,v}(v-u)^2,\\
 J_1&=\sum_{(u,v)}\tau_{u,v}\int_u^v(t-u)t\,dt,
 &J_2&=\sum_{(u,v)}\tau_{u,v}\int_u^v(t-u)t|t|\,dt,\\
 L&=\sum_{(u,v)}\tau_{u,v}\int_{[u,v)}(t-u)\,d\nu(t),
 &D&=(3+\sqrt{10})\rho-\gamma-3J.
\end{align*}
For support in $[-R,R]$, $J\le2R$, $|J_1|\le RJ/2$,
$|J_2|\le R^2J/2$ and $0\le L\le2R$. Thus the sums converge absolutely.
The gap/defect theorem supplies $D\ge0$ when that sign is used below;
the algebraic identity itself does not require that theorem.

Put $\mathcal A=\int\tau(Q_c-h)$. Contact implies
\[
 h(u)-h(v)=\nu([u,v))=Q_c(u)-Q_c(v).
\]
Integrating $h(t)=h(v)+\nu([t,v))$ across the gap gives
\[
 \int_u^v(Q_c-h)
 =\int_u^v(t-u)(-Q'_c(t))\,dt-\int_{[u,v)}(t-u)\,d\nu(t).
\]
The half-open upper endpoint is essential. The slack is zero on the
support and $\tau=0$ outside its convex hull. Substituting $Q'_c$ and
summing gives the first evaluation
\begin{equation}\label{eq:first}
 \mathcal A=\frac{\varphi J}{2\sqrt n}+\frac{G_c}{n^{3/2}}-L,
 \quad G_c=\tfrac32c\beta J+(H\varphi+3c\rho)J_1-3cJ_2.
\end{equation}

Let $X^*$ have density $\tau$, independently of $S_{n-1}$. Signwise
Fubini gives $\E g'(X^*)=\E[Xg(X)]$. Conditional use of that identity
and exchangeability show that $(S_{n-1}+X^*)/\sqrt n$ is the zero-bias
law of $Z=S_n/\sqrt n$. Its CDF at $z$ equals
\[
 T=\E[(Z^2-zZ)\ind_{\{Z\le z\}}],
 \qquad\int\tau(t)h(t)\,dt=T.
\]
The factor $Z^2-zZ$ makes endpoint atoms harmless. Integrating the four
coefficients of \eqref{eq:Q} gives
\begin{equation}\label{eq:second}
 \mathcal A=F-\frac{\gamma\varphi}{2\sqrt n}
 +\frac{H\varphi(3-M_4)/6+cK_q}{n^{3/2}}-T,
 \quad K_q=\frac{M_5}4-\frac{\rho M_4}2-\frac{3\beta\gamma}2+\frac\rho2.
\end{equation}
Multiplying the support contact equality by $X$ and using exchangeability
also gives, for $M=\E[Z\ind_{\{Z\le z\}}]$,
\begin{equation}\label{eq:firstmoment}
 M+\varphi=-\frac{\gamma z\varphi}{2\sqrt n}+\frac{cK_1}{n},
 \quad K_1=\E[X|X|^3]-3\beta-\tfrac32\rho\gamma.
\end{equation}

Define $q_z(x)=(x-z)|x-z|$, and let $Z_0$ be standard normal. Set
\[
 \varepsilon=\E q_z(Z)-\E q_z(Z_0)-\frac{2\gamma\varphi}{3\sqrt n}.
\]
The pointwise hinge identity gives
\[
 T-\Phi(z)=-\varepsilon/2-\frac{\gamma\varphi}{3\sqrt n}
                     +z(M+\varphi)-z^2(F-\Phi(z)).
\]
Substitute \eqref{eq:firstmoment} and
$\sqrt n(F-\Phi)=c\rho+a$ in \eqref{eq:second}. Direct collection gives
\begin{align*}
 \sqrt n\mathcal A&=c\rho+a-\gamma\varphi/6+P_0/n
                                          +\sqrt n\varepsilon/2,\\
 P_0&=H^2(c\rho+a+\gamma\varphi/2)
       -H\{cK_1+\varphi(M_4-3)/6\}+cK_q.
\end{align*}
Equating with \eqref{eq:first} proves the following identity.

\begin{theorem}\label{thm:identity}
Let $U=J_1+(M_4-3)/6$ and define
\begin{align*}
 \mathcal R_c(H)={}&c\{3\beta J/2+3\rho J_1-3J_2-K_q\}
    +H(cK_1+\varphi U)\\
 &-H^2(c\rho+\gamma\varphi/2)+E_n,
 \qquad E_n=-n^{3/2}\varepsilon/2.
\end{align*}
Then, for general $c>0$,
\begin{equation}\label{eq:general}
 (n+H^2)a+nc\rho-\frac{n\varphi}{6}(\gamma+3J)+n^{3/2}L
 =\mathcal R_c(H).
\end{equation}
In particular, at $c=\CE$ and $w=H^2/(2n)$,
\begin{equation}\label{eq:conserved}
 \boxed{(n+H^2)a+nc\rho(1-e^{-w})
   +\frac{n\phi(0)e^{-w}}6D+n^{3/2}L=\mathcal R_c(H).}
\end{equation}
If $a\ge0$ and the physical defect inequality $D\ge0$ holds, all four
terms on the left are nonnegative.
\end{theorem}
The $na$ term comes from the ordinary discrepancy, and $H^2a$ comes from
the quadratic hinge. Replacing their sum by $na$ is false. The formula
uses the raw $H$; in standardized units the coefficient is $n(1+z^2)$.

\subsection{What the Fourier representation still requires}
The unchanged compensated-inversion lemma is the exact retained input
at manuscript Section 5b, lines 143--200. It applies to a finite signed
measure of zero mass and first moment with finite absolute second moment,
for every real raw threshold. Applied to the sum law minus the Gaussian
and cubic signed comparator, it gives
\[
 E_n=\frac{2\sqrt n}{\pi}\int_0^\infty
     \frac{\Im(e^{-itH}V_n(t))}{t^3}\,dt,
 \qquad V_n(t)=f(t)^n-e^{-nt^2/2}(1-in\gamma t^3/6).
\]
The signed comparator has mass one and first moment zero. Compact support
already gives all required moments. No occurrence of the contact
coefficient or the extremal score is present in this inversion step, so
the existing proof applies without modification.
This identity does not bound the integral. Complete low- and
high-frequency envelopes, signed phase uniformity, domain coverage,
moment gates and the infinite tail remain obligations. No nonlattice
decay is available merely from the fixed-target reformulation.

\section{Gaussian weight and phase price}
\begin{corollary}\label{cor:weight}
Suppose $c=\CE$, $a\ge0$, $D\ge0$, $n\ge N\ge2$, and an independently
established raw upper bound gives $\mathcal R_c\le r_0<2N/5$. Then
\[
 (n+H^2)a+\frac{n\phi(0)}6D+n^{3/2}L\le\mathcal R_c.
\]
Any later upper bound $\mathcal R_c\le r$ consequently gives
$e^{-w}\ge1-5r/(2n)$ and $nD\le15000r/997$, using the same directed
constants as the manuscript. If also $nD\le Nd_D$ and
$0.4095\rho-d_D/15>0$, set
\[
 q_{\rm extra}=\frac{1-5r/(2N)}2(0.4095\rho-d_D/15)>0
 \quad\text{when } r<2N/5.
\]
A valid signed envelope
$\mathcal R_c\le A(c,\rho)+F_{\operatorname{sgn}H}(c,\rho)|H|
-Q(c,\rho)H^2$ implies, for $\kappa=997/15000$,
\[
 (n+H^2)a+\kappa nD+n^{3/2}L
 \le A(c,\rho)+F_{\operatorname{sgn}H}(c,\rho)|H|
                       -(Q+q_{\rm extra})H^2.
\]
\end{corollary}
\begin{proof}
Equation~\eqref{eq:general} gives
\[
 \frac{\varphi}{6}(\gamma+3J)
 =(1+H^2/n)a+c\rho+\sqrt nL-\mathcal R_c/n
 >c\rho-2/5>0.
\]
Therefore $c\rho-\phi(0)D/6>0$. Regrouping
\eqref{eq:conserved} leaves the additional nonnegative term
$n(1-e^{-w})(c\rho-\phi(0)D/6)$, which may be dropped. The density and
defect bounds follow by retaining one nonnegative term at a time and
using $c>2/5$, $\phi(0)>997/2500$. Finally
$1-e^{-w}\ge we^{-w}$ and $D\le d_D$ price that same term by
$q_{\rm extra}H^2$. It is subtracted once, not counted twice.
\end{proof}

\section{Transport of a zero-barrier scalar schedule}
There is a useful distinction between the actual score $C_{\rm act}$,
the contact multiplier, and the deleted-sum constants. They coincide or
are compared in the old derivation, but their roles can be separated.
In this note the contact multiplier is always $c$. The predecessors
remain global bounds for actual deleted convolutions. The actual score
is given by \eqref{eq:actual}; it is not set equal to $c$.

\begin{theorem}[Conditional scalar-schedule transport]\label{thm:transport}
Consider a finite, ordered, zero-barrier acquisition and exclusion
schedule at every actual $n\ge N$. Suppose its proof has the following
properties.
\begin{enumerate}
\item Its primitive moment, geometry, Fourier, count and concentration
inequalities hold for every actual law in their stated domains. Its raw
remainder estimates are universal scalar inequalities for the displayed
$\mathcal R_t$, with a positive formal multiplier $t$ in an explicitly
covered interval containing $c$.
\item Its support steps use the zero-loss contact bound; its resource
steps use the conserved identity, licensed weight removal, phase price,
and nonnegativity of the excess. All moment/support domains are acquired
in their stated order for the same law. Every denominator, sign, strict
gap and uniform-index gate is checked.
\item A charged corner bound, if used, subtracts only a nonnegative
charge $\chi(t)\rho$ with $\chi(c)=0$, and its proved scalar envelope
contains the point $t=c$. No positive lower excess barrier is used unless
it has separately been established for this new witness.
\item The confinement proof uses the displayed derivative of the contact
majorant with multiplier $t$ and separate valid deleted-sum bounds. Its
strict budget remains a strict upper envelope for $\mathcal R_c$ after
the phase price.
\item Every positive-noise leaf strictly bounds the count right side
with repair multiplier $t$ on its domain, and its bound at $t=c$ is
negative. The complete cover and the zero-noise strict binomial case
are independently valid.
\end{enumerate}
Then the same scalar bounds and unchanged numerical targets, if proved
on those domains, exclude a maximizer of $A_n(c)$ with $a\ge0$ under
$C_{n-1},C_{n-2}\le c$. An upper bound on $C_{\rm act}$ or on $na/\rho$
is not an entry requirement for this transported schedule.
\end{theorem}
\begin{proof}
Initialize with Theorems~\ref{thm:attain}--\ref{thm:support} and the
universal moment bounds. For a raw scalar inequality specialize its
formal multiplier to $t=c$, without changing the actual marginal or any
of its other parameters. Theorem~\ref{thm:identity} and
Corollary~\ref{cor:weight} supply the required nonnegative resources;
discarding $(n+H^2)a$ only weakens the resulting inequality. At a charged
step the subtracted formal charge is zero. Thus the same endpoint
enclosure dominates the specialized remainder. Support improves by the
extra $-na$, which can be discarded. The explicit confinement derivative
is the derivative of \eqref{eq:Q}, with coefficient $c$, whereas the
concentration estimates use $C_{n-1},C_{n-2}$ independently. A strict
raw envelope remains strict after discarding any of these nonnegative
terms. Induction over the finite ordered schedule establishes every
successive enclosure for this law. Nothing is inferred from an unproved
fixed point.

For the final count step, let $Q$ be the difference from the binomial
CDF and $f(u)$ its normal phase correction, as in the current count
identity. Direct subtraction yields
\begin{equation}\label{eq:transport-master}
 \begin{aligned}
 a+\rho_B(c-C_B)={}&\sqrt n\pi_k\{Q/\pi_k-f(u)\}\\
 &+\sqrt n\left\{\Phi\!\left(\frac{\delta}{\sqrt{nv}}\right)
  -\Phi\!\left(\frac{\delta}{\sqrt{n(v+s)}}\right)\right\}
 +c(\rho_B-\rho).
 \end{aligned}
\end{equation}
The right side is exactly the old comparison right side with formal
multiplier $t=c$. Its strict negative upper bound contradicts $a\ge0$
and $C_B\le c$. The nearest-count mass sign and whole-count curvature
proofs therefore transport unchanged once their right-side inequalities
are proved. At zero noise, maximizing over $H$ for a fixed law still
puts $H$ at a binomial atom, because the penalty is independent of $H$;
the separately assumed strict binomial comparison gives $a<0$.
\end{proof}

\paragraph{Specific scalar bounds that meet the multiplier test.}
The continuous-cell atlas bounds the constant, linear and quadratic
coefficients of $\mathcal R_t$ before using conservation. Its separate
convexity in $t$ and $\rho$ is purely algebraic. Thus its four-corner
bound remains valid at $t=c$ whenever the scalar interval contains $c$;
it does not require $C_{\rm act}$ to lie in that interval.
The two-interval bound in current Section 6 likewise still bounds
$G(c)$: with $C_0<c<\chi<C_1$,
\[
 G(c)\le\max\{G(C_0),G(\chi)\}
 \le\max\{G(C_0),G(\chi),G(C_1)-N(C_1-\chi)\}.
\]
This uses only the old low interval and its convexity, not an upper
bound for $C_{\rm act}$. The coherent six-corner bound also contains
$c<0.40974$, where its charge $n(t-0.40974)_+$ is zero. The fixed-target
calculation can eventually evaluate $G(c)$ directly, but the larger old
bound remains sound if its primitive input bounds have been proved.

The signed cubic repair uses only $0<c_*\le t\le c^*$ as a scalar
multiplier condition. Setting $t=c\in(0.4,0.4226)$ keeps both signs of
the cubic coefficient valid; the negative case still uses its upper
multiplier. No actual-score cap is needed for that version of the repair.
The confinement multiplier interval $(0.4,0.4253]$ also contains $c$;
its two deleted-sum caps instead follow from the first-failure hypothesis.

\paragraph{Limits of this transfer.}
A scalar enclosure originally justified only at a ratio optimizer is
not universal merely because it is stored as numbers. Its derivation
must have the properties in Theorem~\ref{thm:transport}. In particular,
the complete low/high-frequency and floor-42 domains, the initial raw
ceiling licensing weight removal, each strict budget, all acquired
moment intervals, and the complete count cover are still inputs.
Positive-barrier profiles with $K>0$ or coefficient intervals excluding
$c$ are not covered. No assertion here accepts the old numerical
predicates or all retained analytic candidates. This is a proved
conditional transport, not an unconditional LOCAL42 theorem.

\section{A replacement implication, not an established exclusion}
\begin{corollary}\label{cor:global}
Fix $N\ge2$ and $c>0$. Assume $C_j\le c$ for $j<N$. Assume independently
that at every $n\ge N$ with $C_{n-1}\le c$ (and $C_{n-2}\le c$ if an
argument uses that deletion), no maximizer of $A_n(c)$ has $a\ge0$.
Then $C_n<c$ for every $n\ge N$.
\end{corollary}
\begin{proof}
Choose the first failure of the strict conclusion. All required
predecessors are at most $c$. Theorem~\ref{thm:attain} provides the
forbidden excess maximizer with $a\ge0$.
\end{proof}
This new exclusion has no upper-score or upper-excess entry premise, so
the one-step recurrence is unnecessary \emph{for this conditional
implication}. The old LOCAL42 assumption explicitly has both such
premises and concerns a ratio maximizer. It cannot be substituted for
the new assumption. Exclusion only of $a>0$ yields a weak global bound,
not the strict tail.

For a binomial comparator of score $C_B$ and third moment $\rho_B$,
direct subtraction gives
$a=\rho_B(C_B-c)+\mathcal E_c$. The bound
$\mathcal E_c\le\rho_B(c-C_B)$ yields only $a\le0$.
Strict exclusion requires a strict comparison or a separately proved
strict margin. This equality boundary must not disappear during the
reformulation.

\section{Adversarial checks and remaining scope}
The separate exact checker uses sparse Laurent polynomials over rational
numbers to derive the zero-bias moments of $Q_c$, the truncated first
moment, both slack expansions and the hinge elimination. A mutation
using $na$ leaves exactly $H^2a$. Finite-law checks use symmetric and
skew Bernoulli, its reflection, a three-point law with a zero atom, a
four-point law, and three symmetric tiny-far-atom laws. They check 108
zero-bias/hinge instances and 324 raw-mixture derivatives, at negative,
zero, positive and colliding thresholds. Ninety-eight collisions detect
the incorrect closed upper endpoint in the gap measure.

Two further controls identify analytic failure modes. For Bernoulli on
$\{-1,1\}$, the open-event function
$y\mapsto\Pp(X<-y)$ at $y=1$ has value zero and left limit $1/2$, so it
is not upper semicontinuous. For integer $m>1$, put mass $1/(2m^3)$ at
each of $\pm m$, mass $(1-1/m)/2$ at each of $\pm1$, and the remaining
mass $1/m-1/m^3$ at zero. These standardized laws have
$\rho_m=2-1/m$, while the weak limit has $\rho=1$; their third-moment
tails at $m$ equal one. This refutes the unjustified UI3 shortcut, not
the attainment proof above.

\begin{center}
\begin{tabular}{p{43mm}p{48mm}p{49mm}}
\toprule
Item & Mathematical status & Execution status\\
\midrule
Attainment, contact, support transfer & Derived above, at fixed index;
quantitative support keeps its listed inputs & Exact finite-law controls pass\\
Conserved coefficient and sign gates & Exact identities and conditional
consequences above & Independent rational algebra passes\\
Complete Fourier bounds and new local exclusion & Not established by this
note & No native run or certificate acceptance\\
\bottomrule
\end{tabular}
\end{center}
The accompanying Markdown note maps the affected edges to exact current
source locations. No numerical workload saving is certified by this
change of extremal problem alone. Historical certificates remain
untouched, and a bound for a different optimizer cannot inherit their
acceptance status.
